\documentclass{article}
\usepackage{graphicx} % Required for inserting images
\usepackage[margin=1in]{geometry}
\usepackage{parskip}
\usepackage{amssymb}
\usepackage{amsmath}
\usepackage{amsthm}
\usepackage[table]{xcolor}
\usepackage{mathpartir}
\usepackage{stmaryrd}
\usepackage{hhline} 
\usepackage[colorlinks=true, linkcolor=blue, citecolor=blue]{hyperref}

\theoremstyle{definition}
\newtheorem{theorem}{Theorem}
\newtheorem{corollary}{Corollary}
\newtheorem{algorithm}{Algorithm}

\newcommand{\note}[1]{\textcolor{red}{#1}}
\newcommand{\xor}[0]{\textsf{xor}}
\newcommand{\blue}[1]{\textcolor{blue}{#1}}
\newcommand{\red}[1]{\textcolor{red}{#1}}
\definecolor{darkgreen}{rgb}{0,0.5,0}
\newcommand{\green}[1]{\textcolor{darkgreen}{#1}}
\newcommand{\purple}[1]{\textcolor{purple}{#1}}

\title{Automated Reasoning for Network Protocols with Cryptographic Components}
\author{Rob Lorch}
\date{Fall 2026} 

% \renewcommand{\familydefault}{\sfdefault}

\begin{document}

\begin{titlepage}
    \centering

    {\Large \textbf{Thesis Proposal Report} \par}
    \vspace{0.5in}

    {\LARGE \textbf{Automated Reasoning for Network Protocols with Cryptographic Components} \par}
    \vspace{.75in}

    {\large
    Submitted by\par
    \vspace{0.2in}
    \textbf{Rob Lorch}\par
    }
    \vspace{0.55in}

    {\large
    In partial fulfillment of the requirements for the Doctor of Philosophy\par
    Department of Computer Science\par
    The University of Iowa\par
    }
    \vspace{0.75in}

    {\large
    Thesis Committee:\par
    \vspace{0.2in}
    Dr. Omar Chowdhury\par 
    Dr. Katherine Kosaian\par
    Dr. Garrett Morris\par
    Dr. Sriram Pemmaraju\par 
    Dr. Cesare Tinelli\par
    }
    \vspace{1in}

    {\large
    November 2026\par
    }

    \vfill
\end{titlepage}

\section*{Abstract}

Network protocols with cryptographic components are hard to prove correct. 
Defining the notion of protocol correctness is itself nontrivial, 
even before attempting a proof. 
The general difficulty of the problem has resulted in flawed ``proofs" of correctness, 
sometimes leading to real-world security impacts. 
Even worse, security concerns are often eschewed in favor of financial considerations, 
and so proofs are often not attempted in the first place.
In an ideal situation, one would like to propose a security protocol design along with a security property, 
and then have some automated tool either prove the protocol design correct with respect to the property 
or return a concrete attack trace, all with minimal human interaction (we call this goal $G_0$). 
Once the \emph{design} has been proven correct, we would also want to prove the \emph{implementation} correct. 
However, even semi-automated techniques for implementation-level verification are currently 
out of reach; instead, we aspire for automated techniques for discovering implementation-level bugs (we call this goal $G_1$). 
In this thesis proposal, we describe a set of research contributions 
leveraging SMT-based automated reasoning techniques in pursuit 
of goals $G_0$ and $G_1$. 
In particular, we discuss 
$(1)$ a new cryptography-aware model-based protocol testing approach, 
$(2)$ a new proof technique with tool support for cryptographic protocol verification, and 
$(3)$ a new input generation algorithm for rigorous and automated implementation-level testing.
We outline and evaluate these contributions, 
discuss completed and proposed work, 
and suggest a possible timeline to complete the remaining tasks.

\section{Introduction}




Security protocols are hard to get right. Examples abound of the discovery of inherent, 
design-level issues in security protocols only \textit{after} said protocols were designed, implemented, 
and in use \cite{lteinspector,5greasoner,ns_attack}. 
In this document, we argue that approaches based on \emph{formal methods} can help address these issues 
by aiding in the detection of security flaws at the project design level.
An often-cited example is the infamous Needham-Schroeder mutual authentication protocol \cite{needham_schroeder}, 
which we discuss in Section \ref{sec:sym-attacks}; its flaw was only discovered 17 years after the protocol 
was proposed with the help of formal tools \cite{ns_attack}.
Of course, there are also modern examples. A particularly relevant but tricky example 
is the 4G LTE/5G protocol stack \cite{4g_lte,5g}. 
The protocols are composed of various phases and layers, and they include the composition of cryptographic reasoning 
with non-cryptographic reasoning (e.g., linear integer arithmetic). 
Formal approaches to analyzing the 4G LTE and 5G protocol stacks have shown promise: 
they have resulted in the discovery of numerous design-level flaws, 
even when making the overly optimistic simplifying assumption that all cryptographic primitives 
are perfectly secure \cite{lteinspector,5greasoner,dolev_yao}. 
However, the state-of-the-art techniques \cite{lteinspector,5greasoner} still rely on manual interaction 
that requires expertise in \emph{two} separate formalisms, potentially suggesting an area for future improvement.

Late discovery of design-level problems has huge costs, financial and otherwise---it is estimated 
that problems discovered earlier at the \textit{design stage} (rather than the \textit{implementation stage}) 
of a software project
generally cost 5-30 times \emph{less} to address \cite{fm}. 
The current industry-standard method to detect software flaws is \textit{testing}. 
However, testing has limitations: $(i)$ it relies on the existence of a concrete implementation, 
$(ii)$ it can only prove the presence of bugs and never their absence, 
and $(iii)$ designing test cases that account for the wide range of possible adversarial influence is difficult. 
% {\begin{itemize}
%     \item Give motivation for reasoning about cryptography; Needham-Schroeder protocol
%     \item From Omar -- look at LTEInspector/5GReasoner work for motivation; real-world protocols have cryptographic components; one can forget about crypto but get spurious counterexamples; otherwise one has to integrate cryptographic reasoning. Even in cases where primitives are assumed secure there are real-world attacks. There are current approaches, model-based testing and abstraction/refinement to find these attacks. But they have downsides. But we need to take a step back and build a story based on details.
% \end{itemize}}
Therefore, to facilitate early detection of design flaws, we are interested in the formal verification of 
\textit{security protocols} with respect to \textit{security properties}, 
even under the influence of an \textit{adversary}, 
dubbed the \textit{security protocol verification problem} 
(where the notions of \textit{security protocol}, \textit{security property}, and \textit{adversary} will be elaborated later). 
We are particularly interested in \emph{(semi-)automated approaches} to addressing the 
security protocol verification problem, keeping in mind two core accompanying concerns,
\emph{expressivity} and \emph{tractability}:
approaches must be \emph{expressive enough} to model modern protocol features such as cryptographic primitives with 
algebraic properties (e.g. exclusive or), arithmetic constraints, and non-trivial control flow, 
while still remaining \emph{tractable enough} to produce proofs or counterexamples with minimal human interaction. 

\dots  

The SMT revolution has had far-reaching consequences in general-purpose 
software verification (eg Dafny), model checking (eg general-purpose model checkers), 
and bug finding (symbolic execution). 
However, its impact has been limited in the context of 
cryptographic systems. 
Reasoning rigorously about mathematical cryptography is traditionally outside the scope of 
automated approaches (ofc there have been some successes like Cryptol); 
reasoning wrt an abstraction has relied on things that SMT solvers are historically weak at (ADTs and quantifiers); 
implementation-level reasoning in the fuzzing world has traditionally viewed constraint solving as too expensive 
to assist in the generation of test cases.

\textbf{Roadmap.} Roadmap here.

\section{Background}
Network protocol, security protocolf

Design vs implementation level.

Symbolic vs computational model.

Blackbox vs graybox vs whitebox fuzzing, 
motivation for each, and generative vs mutation-based algos.


\section{Thesis Statement \& Contributions} 
\emph{This thesis leverages SMT-based automated reasoning techniques to advance
the state of the art in the analysis of network protocols with cryptographic 
components.} We focus on two key goals. 

\noindent\textbf{Goal $G_0$.}
We are interested in \emph{semi-automated}, \emph{design-level verification}
of these protocols.


\noindent\textbf{Goal $G_1$.}
We are interested in \emph{fully automated}, \emph{implementation-level bug finding}. 

To address these two key goals, we propose three research contributions.

\noindent\textbf{Contribution $1$.} 
To address $G_0$, we develop a cryptography-aware model-based testing approach 
that offers higher degrees of automation and expressivity compared to prior work. 
We apply this approach to the Uptane \cite{} over-the-air update protocol, exhibiting X 
previously unknown design-level security vulnerabilities. 

\noindent\textbf{Contribution $2$.}
To further address $G_0$, 
we note a fundamental weakness of the model-based testing approach presented in Contribution $1$: 
it requires two separate formal models, 
one in a general-purpose model checker and one in a dedicated cryptographic 
protocol verifier (CPV). 
Further, we note a fundamental weakness of prior CPVs: 
the proofs from the standard CPVs \cite{tamarin,proverif} are inherently 
monolithic, and hence verification failures are opaque and thus difficult for users to address. 
Aside from these standard approaches, 
there also exist compositional techniques that offer more information in the case of verification failure 
\cite{dystar}, but their proof methodologies offer little automation. 
To address these weaknesses, we present a new cryptographic protocol verification approach 
based on deductive verification that only requires a single formal model, 
and additionally advances prior compositional proof techniques to increase automation potential
(compared to existing CPVs).
We demonstrate the utility of this approach by presenting proofs 
of Y $\langle$protocol, property$\rangle$ pairs with higher levels of automation than were possible in prior work. 

\noindent\textbf{Contribution $3$.}
Despite the advances of Contributions $1$ and $2$, we note that there are certain classes of 
bugs that are impossible to address with these approaches---namely, bugs that only exist 
in concrete protocol implementations, and not at the protocol design phase. 
Therefore, to address $G_1$, we generally resort to implementation-level fuzz testing. 
Prior works on fuzzing these protocols have shown promise with finding vulnerabilities in practice, 
but they suffer from fundamental weaknesses in $(a)$ requiring access to the system under tests's (SUT's)
source code or binary, and in $(b)$ often failing to produce inputs that drive deep enough 
into the application's logic to reliably trigger bugs.
To address these weaknesses, we develop an SMT-based input generation approach 
to fully automatically generate test cases that conform to an input specification 
phrased as a context-sensitive grammar.
We demonstrate the utility of the approach by measuring the efficiency of the input generation 
algorithm and the quality of the generated inputs. 
We additionally demonstrate Z bugs on the WiFi SAE handshake protocol that were 
discovered by our approach and out of reach by prior approaches \cite{saecred}.

\section {Contribution 1}
Weaknesses of prior work
\begin{itemize}
  \item Borrow from Uptane paper. Recreate fig of lazy approach. Also endless attack chaining is it possible?
  \item We want to reason about systems with cryptographic components. Historically, this has been split 
  into reasoning about \emph{cryptographic protocols} and \emph{(cryptographically-abstracted) network protocols}. 
  But in reality these things mix. Cite LTEInspector/5GReasoner.
\end{itemize}

\section {Contribution 2}
Weaknesses of prior work
\begin{itemize}
    \item 
  Following up on Uptane work,
  we would like to do everything in one model, and what if the threat model cannot be adequately 
  summarized by the proposed standard/replay/bounded replay/no attackers?
  \item See SP submission challenges of ProVerif and Tamarin, and also challenges of invariant-based 
  approaches with Figure. Maybe also figure of Tamarin proof in progress.

  And compression theorem is protocol-specific
  \item Maybe use big CPV table from SP submission
\end{itemize}


\section {Contribution 3}
Weaknesses of prior work
\begin{itemize}
  \item This is great, but we are still missing some bugs. 
  Lots of bugs only exhibit themselves at the implementation level; recall these type confusion attacks.
  \item See ICSE '26 and perhaps the slides for summary and diff with prior work (perhaps that table). 
  Fandango has problems with non-monotonic constraints; ISLa is a string solver and does not leverage 
  incrementality of SMT solvers. Perhaps give bogus ISLa semantics rule quantifying over integers 
  as an example. 
  \item Also look into QuickCheck \& Co, and SyGuS engines, CLP, and so on.
  \item Basic challenges of fuzzing network protocols with cryptographic components, especially in black box. 
  We focus on black box because in whitebox and graybox there are pretty well-established techniques 
  (and it turns out we can still beat graybox fuzzers)
\end{itemize}


\section{Completed \& Proposed Work}
\subsection{Completed Work}
\begin{itemize}
  \item Uptane (RAID '24), number of vulnerabilities
  \item Goblin (ICSE '26, S\&P '25), number of vulnerabilities
  \item Medjed submission (S\&P '27), results overview 
\end{itemize}

\subsection{Proposed Work}

Goals (have not distinguished between main and stretch yet)
\begin{itemize}
  \item Proper input language (give examples where Medjed language is not really desirable)
  \item Multiple reasoning engines (Kind 2, Tamarin for cex to induction, compilation to Tamarin)
  \item Find counterexamples 
  \item Dedicated neuro-symbolic reasoning 
  \item Broader eval on benchmarks 
  \item More novel protocol analysis problems
\end{itemize}

\section{Timeline}
reasonable and detailed; flexibility (perhaps flex time?); >= 2 months for dissertation writing; give time for comments and responding to feedback before official deadline for thesis submission; go backwards; Omar did 3 months for revisions
Question: Was Omar's 3 months after the defense and before the final submission, or before the defense?
\begin{itemize}
  \item Now till June 2027: Work on lemma synthesis paper; revise S\&P '27 paper if rejected
  \begin{itemize}
    \item Lemma synthesis algorithm 
    \item Tool support 
    \item Evaluation
    \item Paper writing
  \end{itemize}
  \item June 2027: Submit lemma synthesis paper to S\&P
  \item July 8, 2027: Start writing thesis
  \item Early September 2027: Completed draft of thesis for committee
  \item Early October 2027: Thesis defense
  \item December 2027: Thesis deposit due
\end{itemize}

\section{Notes}

\begin{itemize}
  \item At least 15 pages, not beyond 30 
  
  \item Introduction and thesis statement
  \item Enough background to make concepts clear
  \item Shortcomings of current work
  \item Contributions of the thesis
  \item Completed work
  \item Proposed work
  \item Timeline and plan of attack
  \item Contingencies (plan B, C, etc.)
  \item Think in terms of achievements, not papers
  \item How I intend to go about the remaining work 
  with \emph{some way to measure success}
  \item First draft with full ambition
\end{itemize}

\bibliographystyle{plain}
\bibliography{biblio}
\end{document}

